\begin{thema}{A}
\begin{erwthma}
    \item Να αποδείξετε ότι η παράγωγος της συνάρτησης $ f(x)\cdot g(x)\cdot h(x) $ του γινομένου τριών παραγωγίσιμων συναρτήσεων ισούται με \[ [f(x)\cdot g(x)\cdot h(x)]'=f'(x)\cdot g(x)\cdot h(x)+f(x)\cdot g'(x)\cdot h(x)+f(x)\cdot g(x)\cdot h'(x) \]
    \item Πότε μια συνάρτηση $f$ ονομάζεται συνεχής; Πότε συνεχής σε κλειστό διάστημα $[a, \beta]$;
    \item Τι ονομάζουμε αρχική συνάρτηση ή παράγουσα μιας συνάρτησης $f$ σε ένα διάστημα $\Delta$;
    \item Να χαρακτηρίσετε τις προτάσεις που ακολουθούν, γράφοντας στο τετράδιό σας, δίπλα στο γράμμα που αντιστοιχεί σε κάθε πρόταση, τη λέξη \textbf{Σωστό}, αν η πρόταση είναι σωστή, ή \textbf{Λάθος}, αν η πρόταση είναι λανθασμένη.
    \begin{alist}
        \item Αν $\dint_a^\beta f(x) dx = 0$ και η $f$ παίρνει και θετικές και αρνητικές τιμές, το εμβαδόν του αντίστοιχου χωρίου είναι στο σύνολο 0. (Το εμβαδόν είναι $>0$).
        \item Αν υπάρχει εξακριβωμένα $\xi \in (a, \beta)$ με $f'(\xi)=0$, τότε υποχρεωτικά $f(a)=f(\beta)$.
        \item Η συνάρτηση $f(x) = \sqrt{x}$ είναι παραγωγίσιμη στο $x=0$.
        \item Το πεδίο ορισμού του αθροίσματος $f+g$ είναι πάντοτε η τομή των πεδίων ορισμού τους.
        \item Αν $\dint_a^\beta f(x) dx = \dint_a^\beta g(x) dx$, τότε υποχρεωτικά συνεπάγεται ότι $f(x) = g(x)$ σε όλο το $[a, \beta]$.
    \end{alist}
    \item Αν η παράγωγος $f'(x_0)$ ΔΕΝ υπάρχει (ενώ η $f$ είναι συνεχής στο $x_0$), τότε στο $x_0$:
    \begin{tasks}[label=\let\textdexiakeraia\relax\alph*.](2)
        \task Η $f$ δεν μπορεί να έχει ακρότατο.
        \task Η $f$ μπορεί να έχει ακρότατο (π.χ. $f(x)=|x|$ στο $0$).
        \task Η $f$ είναι αναγκαστικά σταθερή κοντά στο $x_0$.
        \task Η $f$ δεν ορίζεται στο $x_0$.
    \end{tasks}
\end{erwthma}
\end{thema}
